% Introduction and Related Work sections drafted by literature-agent.
% These are intended to be incorporated into main.tex directly (no \input).

\section{Introduction}

Protein function lies at the heart of biology, medicine, and pharmacology, yet the experimental methods that assign function to individual proteins remain slow and expensive. High-throughput sequencing has outpaced annotation, leaving the majority of newly deposited proteins uncharacterized. Computational comparison is the principal lever for closing this gap: given one or two proteins whose roles are known, one can infer the role of a third by asking how similar it is to the first two. Two broad traditions dominate this inference. Sequence alignment, exemplified by PSI-BLAST and its successors \citep{boratyn2012delta}, compares amino-acid strings using substitution matrices and evolutionary profiles. Structural alignment, embodied by RMSD, TM-Score, TM-align \citep{zhang2005tmalign}, and graphlet-based tools such as GR-Align \citep{maloddognin2014gralign}, compares three-dimensional coordinates and is more sensitive at remote evolutionary distances.

Recent advances have pushed both traditions forward. AlphaFold2 \citep{jumper2021alphafold} and the language-model approach ESMFold \citep{lin2023esmfold} have removed structure prediction as a bottleneck, and OmegaFold \citep{wu2022omegafold} demonstrates accurate prediction from single sequences. On the comparison side, TM-Vec \citep{hamamsy2024tmvec} maps structures into a vector space using a transformer trained with TM-score supervision, enabling rapid homology search. Nevertheless, alignment-based pipelines remain computationally expensive at proteome scale, while deep-learning embeddings are opaque: their numerical dimensions do not correspond to physical quantities, and their accuracy on rare folds is constrained by training data.

We introduce an alignment-free descriptor that is simultaneously fast, accurate, and physically interpretable. Statistical energy functions derived from known protein structures empirically capture the most probable state of a protein and describe the microstates of interactions within the native fold \citep{koppensteiner1998knowledge,sippl1990potential}. The prevailing assumption is that a native protein structure occupies a minimum-energy state: the closer a structure is to native, the lower its total energy. We take a step beyond this assumption and hypothesize that two proteins that are structurally or evolutionarily similar possess analogous \emph{energy profiles}. Concretely, each protein is represented by a 210-dimensional vector whose entries sum pairwise amino-acid interaction energies computed from a distance-dependent knowledge-based potential. This Structural Profile of Energy (SPE) is computed directly from atomic coordinates. We further show that a corresponding Compositional Profile of Energy (CPE) can be estimated from amino-acid composition alone, eliminating the requirement for a solved structure.

The idea of using energy as a compatibility signal is not new. Bowie, L\"uthy and Eisenberg \citep{bowie1991eisenberg} pioneered the notion of mapping amino-acid sequences to structural environments using an environmental profile of energy, and Bryngelson, Onuchic and Wolynes \citep{bryngelson1995wolynes} formalized the energy-landscape theory of folding with a spin-glass Hamiltonian. Dosztanyi and colleagues \citep{dosztanyi2008energy} estimated energy from amino-acid composition. Our contribution is to collapse such energetic information into a fixed-length descriptor that supports pairwise dissimilarity computation via the Manhattan distance, with no alignment step required. The Manhattan distance is interpretable (it decomposes over the 210 interaction types) and computes in time linear in the descriptor length.

We validate the framework on six settings of increasing difficulty: (i) SCOP class/fold/superfamily/family hierarchy on ASTRAL40 and ASTRAL95 domains \citep{chandonia2019scope}; (ii) head-to-head 1-NN classification on a 251-domain CATH benchmark against GR-Align, RMSD, TM-Score, Yau-Hausdorff distance, and TM-Vec; (iii) phylogeny of the ferritin-like superfamily \citep{lundin2012ferritin} in the twilight zone of sequence similarity; (iv) clustering of 143 coronavirus spike glycoproteins curated from CoV3D \citep{gowthaman2021cov3d}; (v) classification of 690 bacteriocins from the BAGEL database \citep{alkhalili2018bagel} together with Papain-like protease domains across four Betacoronavirus subgenera; and (vi) a drug-combination application that correlates a new energy-profile separation measure with the network-based score of Cheng and colleagues \citep{zhou2020network}. On the CATH benchmark, CPE attains 100\% accuracy in 1 second compared to 67 seconds for TM-Vec and classical baselines between 59.2\% and 81.5\%. On the coronavirus spike task, CPE reaches an Adjusted Rand Index of 0.95 in 0.9 seconds, with SPE achieving a perfect ARI of 1; TM-Score requires 9.7 hours.

Our contributions are as follows. We formalize a sequence-derived energetic descriptor and show that it correlates strongly with its structural counterpart across ASTRAL40 and ASTRAL95 with p-values below $10^{-16}$. We demonstrate that the descriptor captures hierarchical, functional, and evolutionary signal, including beyond the twilight zone where sequence alignment alone fails \citep{chothia1986twilight}. We benchmark against classical and deep-learning baselines and establish a Pareto frontier that is both more accurate and orders of magnitude faster. Finally, we extend the method to drug-combination scoring and proteome-scale screening of 4,405 coronavirus models across 28 families.

\section{Related Work}

\paragraph{Knowledge-based statistical potentials.}
Statistical potentials infer effective pairwise energies from the observed frequencies of atom--atom contacts in high-resolution structures. Sippl's seminal formulation \citep{sippl1990potential} and later refinements \citep{koppensteiner1998knowledge} established distance-dependent potentials as a workhorse for threading, fold recognition, and model quality assessment. We adopt a distance-dependent variant of the Sippl potential as the foundation of our descriptor.

\paragraph{Energy landscape theory.}
Energy-landscape ideas, culminating in the Bryngelson--Wolynes funnel picture \citep{bryngelson1995wolynes}, cast the native fold as a global energy minimum reached through minimally frustrated pathways. Recent work has exploited residual frustration to probe functional divergence in homologs such as $\alpha$- and $\beta$-globins \citep{freiberger2023frustration}, motivating our use of pairwise energy distributions to distinguish close subfamilies.

\paragraph{Sequence-to-fold mapping via energy profiles.}
Bowie, L\"uthy and Eisenberg \citep{bowie1991eisenberg} introduced an environmental profile that scored the compatibility of an amino-acid sequence with a structural template. Dosztanyi et al. \citep{dosztanyi2008energy} showed that the total stabilization energy of a protein can be estimated directly from sequence composition. Our CPE generalizes this line of work to the 210-dimensional space of amino-acid pair types, yielding a vector descriptor rather than a scalar quality score.

\paragraph{Structural comparison metrics.}
TM-align \citep{zhang2005tmalign} combines dynamic programming with a scoring function that normalizes against protein length; its TM-Score has become the de-facto standard for reporting structural similarity. GR-Align \citep{maloddognin2014gralign} extracts topological graphlets from residue contact graphs for flexible alignment. Root-mean-square deviation, computed after least-squares superposition, is simple and widely used but sensitive to outlier regions and flexible loops. Hausdorff-type distances treat structures as point sets and measure maximal bidirectional separation. These methods differ in their balance of accuracy, speed, and robustness but share a dependence on either alignment or explicit atomic coordinates.

\paragraph{Deep-learning embeddings.}
TM-Vec \citep{hamamsy2024tmvec} encodes structures into a fixed-length vector space using a transformer trained on pairs of structures with their TM-scores as supervision; cosine similarity in the embedding approximates TM-Score. The approach achieves strong performance on remote homology detection but inherits two classical limitations of deep learning: reliance on a training distribution and limited interpretability. Our energy profile complements such embeddings by offering dimensions that correspond directly to biophysical interaction types.

\paragraph{Structure prediction.}
AlphaFold2 \citep{jumper2021alphafold}, OmegaFold \citep{wu2022omegafold}, and ESMFold \citep{lin2023esmfold} have transformed the availability of protein structure models, but rely on ensembles of predictors or language models with large compute footprints at inference. CPE trades explicit 3D modeling for a lightweight composition-based estimate that correlates tightly with the SPE derived from atomic coordinates.

\paragraph{Protein classification and phylogeny.}
Hierarchical classifications such as SCOP \citep{murzin1995scop}, CATH \citep{orengo1997cath}, and SCOPe \citep{chandonia2019scope} organize the protein universe at class, fold, superfamily, and family levels. In regions where sequence similarity falls below 20--35\%---the so-called twilight zone \citep{chothia1986twilight}---phylogenetic reconstruction requires structural signal. Lundin and colleagues \citep{lundin2012ferritin} constructed structural evolutionary networks for the ferritin-like superfamily using distance-based NeighborNet methods implemented in phangorn \citep{schliep2011phangorn} and SplitsTree \citep{huson2024splitstree}. We use these same datasets to evaluate whether an energy-based metric recovers the manually annotated topology.

\paragraph{Coronavirus structural biology.}
CoV3D \citep{gowthaman2021cov3d} curates high-resolution coronavirus protein structures and their antibody complexes, and has become a standard resource for comparative analyses of the spike glycoprotein and its receptor-binding domain. MAFFT \citep{rozewicki2021mafft} supplies the multiple-sequence alignment baseline for this domain.

\paragraph{Network pharmacology.}
Cheng and colleagues \citep{zhou2020network} introduced a network-based separation measure $S_{AB}$ between pairs of drug targets in the human interactome. Drug pairs with $S_{AB} < 0$ share a network neighborhood; pairs with $S_{AB} > 0$ are topologically separated. We define an analogous energy-based separation measure $E_{AB}$ on the drug targets' energy profiles and show that the two measures are strongly correlated on a 65-drug antihypertensive panel. In contrast to the network score, $E_{AB}$ requires only protein sequences.

\paragraph{Training set curation.}
Following best practice for knowledge-based potentials, we curate the training set of non-redundant protein chains using PISCES \citep{wang2003pisces}, with atomic contacts identified via Delaunay tessellation implemented in Qhull \citep{barber1996qhull}.
